\section{Limitations and Conclusion}
\label{sec:limitations}

\paragraph{The models are fitted on a moving target.}
Mean lead time falls 5.34 days between training and test, and a model fitted on
the slower past is systematically off-level on the faster present ---
Section~\ref{sec:results-segments} measures that as a positive bias in every
region. Hence $R^2$ against a test-period variance is negative even for a useful
model, which is why MAE leads, and hence the models need refitting on a rolling
basis rather than once \citep{gama2014drift}: a deployment should refit monthly
and monitor the gap between predicted and realised lead time as the trigger.
Calibration belongs in the same cycle --- the probabilities hold on the test
window, but that is a two-month-ahead check, and a calibration fitted on one
prior drifts as the prior does.

\paragraph{The largest driver of lateness is not observable at prediction time.}
The late rate swings between 1.4\% and 21.4\% month to month with Black Friday
2017 and the May 2018 truckers' strike. These are exogenous shocks --- carrier
capacity, national road blockades --- that no order-level feature available at
checkout can anticipate. This, not the choice of algorithm, is the binding
constraint on Problem 1B, and it is why the ceiling was stated in advance rather
than a weak classifier presented as a disappointment.

\paragraph{What the data does not contain, and what rests on assumption.}
The service-quality literature identifies product expectation gaps, seller
communication and returns handling among the drivers of dissatisfaction
\citep{parasuraman2005esqual, rao2011elsq}; none appears in these nine tables, so
the Problem 2 residual is unobserved cause rather than noise, and the free-text
review columns (present for 11.7\% and 41.3\% of reviews) are too sparse to
recover it. Olist publishes no cost figures, so Section~\ref{sec:cost} states its
assumptions and reports break-even precision --- the quantity that decides the
recommendation and moves with the assumption. And selection votes on three
out-of-sample windows covering roughly February to June 2018: the right trade
under drift, but a smaller sample than cross-validation uses, and
Section~\ref{sec:results-late} records a case where the rule preferred a model
that does not transfer. We report that rather than quietly selecting on the test
window.

\paragraph{Generalisation, and what these scores are not.}
Everything here is Brazil, 2016--2018, on one marketplace with a distinctive
geography: 59.7\% of sellers sit in São Paulo state against 42.0\% of customers,
so most orders cross the country and the median one travels 434\,km --- the
distance-driven structure the models learn is a property of that asymmetry. The
outputs are triage scores for internal prioritisation, not entitlements, and must
not be used to penalise a seller, price a customer, or drive an automated
decision about someone without review. The review-risk model in particular is
largely a restatement of how late a parcel was, and a seller penalised by it
would in many cases be penalised for a carrier's failure.

\subsection*{Conclusion}

Two problems, at two prediction moments, from one dataset and one feature
pipeline. The regression reaches MAE 3.90 days --- 41\% below a constant fitted
on the past and 60\% below the platform's own promise --- and its output,
thresholded against that promise, makes a better late-order flag (2.89$\times$
the no-skill floor) than any classifier trained directly on the binary label
(1.24$\times$), so the brief's dual framing turns out to be the substantively
better design. The satisfaction model reaches 3.75$\times$ its floor on a target
where always predicting ``satisfied'' scores 89\% accuracy and catches nothing,
and scores identically on customers it has never seen.

The methodological findings travel further. Cross-validation inside a drifting
training period reorders the candidates relative to the forward window, and a
two-month validation period is short enough to reorder them again --- so
selection belongs on several recent out-of-sample windows judged by rank, and
even then it can fail. Weighting the loss for class imbalance leaves ranking
unchanged and inflates predicted risk threefold, so the threshold is the right
lever. Raw PR-AUC is not comparable across periods with different base rates
while its lift over the floor is. And the most useful business result is a
negative one: a flag with ROC-AUC 0.782 is genuinely informative and still does
not justify an expensive intervention, because its precision never clears the
break-even the cost ratio demands.
